% /root/monodromy, 2026-11-19. Complete proof; independent parent audit passed.
\subsection{A finite collision reduction for the degree-three norm family}
\label{subsec:ss-level3-collision-reduction}
\status{Proved and independently audited. The parent and reduction-descent agents read and checked the
complete reduction, including the conductor stable-rank formula, the
local-ring Frobenius-unit graph comparison and every incidence and
component-assignment step. All three global witnesses have complete
independent proof and program audits. The later
Theorem~\ref{thm:ss-all-odd-full-groups} gives the full image in
every odd degree, using the subsequent uniform collision proofs.}

Work over an algebraically closed field of characteristic two. Let
$E:y^2+y=x^3$, $X=E^2$, $L=\mathcal O(3O)\boxtimes\mathcal O(3O)$,
and choose $\omega\in\mathbf F_4$ with $\omega^2+\omega+1=0$.
Put $i(x,y)=(x+1,y+x+\omega)$, $b=1+i$, and
\[
 \phi=\begin{pmatrix}1&b\\-b^\dagger&1\end{pmatrix}:X\to X.
\]
Then $b^\dagger b=2$, $\phi^\dagger\phi=3$, and
$K=\ker\phi=\{(-bP,P):P\in E[3]\}$ has order nine. The norm gives
\[
 \nu:\mathbf P(H^0(L))=\mathbf P^8\longrightarrow |L^3|,
 \qquad\nu^*\mathcal O(1)=\mathcal O(9).
\]
There are no base points in this map: the norm of a nonzero section
is nonzero. Its exact target line bundle was proved in
Theorem~\ref{thm:ss-uniform-GL8-block}.
In the monomial tensor basis put
$D(s)=\det(s_{ij})$ and $Q(s)=s^{[2]t}(J_3\otimes J_3)s$,
where $J_3$ reverses the three basis vectors. The source Hasse
polynomial is $D^2Q$.

\begin{lemma}[The finite kernel has three collision types]
\label{lem:ss-level3-collision-orbits}
A group of product automorphisms preserving $K$, together with
translations by $K$, acts transitively on nonzero kernel differences.
Its unordered triples have precisely two orbits: a three-point affine
line in $K\simeq\mathbf F_3^2$, and a noncollinear triple.
One may use the representatives
\[
 G=((\omega,\omega^2),(0,0)),\qquad
 H=((\omega^2,\omega^2),(1,\omega)),
\]
with pair difference $G$ and triples $\{0,G,-G\}$ and $\{0,G,H\}$.
Every nonzero kernel difference has both elliptic coordinates nonzero.
\end{lemma}
\begin{proof}
The automorphisms $i$, $j(x,y)=(x+\omega,y+\omega^2x+\omega)$,
and $z(x,y)=(\omega x,y)$ preserve $O$ and the cubic. Their action
on $E[3]$ generates $\mathrm{SL}_2(\mathbf F_3)$.
This is also checked without any group identification in the retained
finite certificate: the nine points of $E(\mathbf F_4)$ are all of
$E[3]$; the displayed generators give twenty-four distinct permutations,
all preserving addition, and their action on nonzero points is transitive.
The action is faithful, since an automorphism of the cubic fixing all
nine flexes is the identity; it preserves the Weil pairing, so its
order identifies it with $\mathrm{SL}_2(\mathbf F_3)$.

The map $-b$ on $E[3]$ is invertible and has determinant two. For
$v\in\mathrm{SL}_2(\mathbf F_3)$ the conjugate
$u=(-b)v(-b)^{-1}$ again lies in $\mathrm{SL}_2(\mathbf F_3)$.
The product automorphism $(u,v)$ preserves the graph $K$.
Its action on the graph is the action of $v$ on the second coordinate.
The resulting affine special linear action has order $9\cdot24=216$.
It is transitive on nonzero differences. On unordered triples, the
collinear triples form one orbit of size twelve. A noncollinear triple
can be translated to contain zero, and a reordering changes the sign
of its nonzero determinant, so all such triples form the other orbit,
of size seventy-two. Since $b$ is invertible on $E[3]$, neither
coordinate of a nonzero graph point is zero. Direct addition gives
the stated $G,H$.

The retained exact program
\path{agents/monodromy/code/supersingular_level3_kernel_orbits.py}
checks the full addition table, the twenty-four automorphisms, conjugation
closure, all graph points and both affine orbits. Its output is
\path{agents/monodromy/outputs/supersingular_level3_kernel_orbits.json};
the author run passed in $0.007294220002222573$ seconds.
\end{proof}

For $g\ne0$, let $R_g$ denote the closure of the pair-tangency incidence
image: $C_s$ contains $P,P+g$ and the two translated tangents coincide.
For distinct $0,g,h$, let $T_{g,h}$ denote the closure of the locus
where $C_s$ contains $P,P+g,P+h$.

\begin{lemma}[Irreducibility and separation from the two source factors]
\label{lem:ss-level3-collision-incidence}
Each $R_g$ has one possible divisorial component arising from its
main incidence, and no additional divisorial component from the
exceptional jet locus. Each $T_{g,h}$ is an irreducible incidence
image. If these images contain a smooth noninvariant curve, they
are irreducible divisors. None of these divisors equals $D=0$ or
$Q=0$.
\end{lemma}
\begin{proof}
For fixed $g$, parametrize a common tangent by $(P,\ell)$ with
$P\in X$ and $\ell\in\mathbf P(T_PX)$. Containment and vanishing
of the derivative along $\ell$ at both translated points give four
linear conditions on $H^0(L)$. Since both coordinates of $g$ are
nonzero, the tensor jet calculation in
Theorem~\ref{thm:ss-uniform-GL8-block} gives rank six for the joint
first jets except at finitely many $P$, and rank at least five there.
Thus the four imposed conditions have rank four on the main open
and at least three at those exceptions. The main incidence is a
$\mathbf P^4$-bundle over an irreducible three-dimensional base,
so is irreducible of dimension seven. The exceptional part has
dimension at most $1+(8-3)=6$.

For triples, the three value tensors are independent everywhere.
Indeed three distinct rank-one Segre tensors can be collinear only
if one factor is constant, whereas both projections of $K$ are
injective. The incidence is therefore a $\mathbf P^5$-bundle over
$X$, irreducible of dimension seven.
On a smooth curve not invariant under a nonzero kernel translation,
all translated intersections are finite. Hence the incidence maps
have finite fibres there. The invariant-curve locus cannot dominate:
a nonzero three-torsion translation acts with three distinct eigenvalues
on $H^0(E,\mathcal O(3O))$, since a repeated projective eigenspace
would be a fixed line intersecting the cubic. Its tensor action has
eigenspaces of dimension three, so invariant sections form only a
union of $\mathbf P^2$'s. This proves the divisor assertions under
the stated nonemptiness hypothesis.

No vector subspace of dimension greater than four lies in $Q=0$.
If a subspace $U$ did, polarization of the cubic would imply
$U^{[2]t}JU=0$, so $U\subset(JU^{[2]})^\perp$ and
$2\dim U\le9$. The generic coefficient fibres in the two incidences
have vector dimensions five and six. Thus neither incidence image
can be the cubic $Q$.

For $D$, normalize the two point-value directions in each factor to
$e_0,e_1$. At a generic translated pair the first-factor tangent at
either point has nonzero $e_2$ coefficient. Choose a common tangent
with nonzero first-factor component. The two value equations are
$S_{00}=S_{11}=0$. Set $S_{01}=S_{10}=1$; the two tangent equations
solve for $S_{20},S_{21}$ after the other entries are chosen.
Neither tangent equation involves $S_{22}$, and its coefficient
in $\det S$ is one. Thus a generic pair-tangent coefficient fibre
contains an invertible matrix.

For a triple, its first-factor value directions are generically
independent: dependence would say $3P_1+g_1+h_1=0$ on the elliptic
curve, a finite condition on that coordinate. Normalize them to a
basis. The three rows of $S$ must lie respectively in the three
distinct hyperplanes orthogonal to the second-factor values.
One can choose three independent row vectors in these planes:
otherwise their sum would have dimension at most two and all
three planes would coincide. This gives an invertible $S$ in the
coefficient fibre. Hence neither incidence image equals $D$.
\end{proof}

\begin{proposition}[A reduction to three global geometric witnesses]
\label{prop:ss-level3-three-witness-reduction}
Suppose the following three witnesses exist.
\begin{enumerate}
\item A smooth source curve on $Q\setminus D$ whose norm has one
simple self-tacnode of difference $G$ and otherwise only ordinary
nodes, with unequal values of its extra Cartier-killed form on
the two tangent branches.
\item A smooth source curve with $D Q\ne0$ whose norm has one
three-branch singularity of the tangent-triple type in
Lemma~\ref{lem:ss-tangent-triple-bridge}, with source
triple in the affine-line orbit, and otherwise only ordinary nodes.
\item The same kind of sharp-source tangent-triple witness for the
noncollinear orbit.
\end{enumerate}
Then the degree-three relative Hasse divisor on $|L^3|$ is
geometrically irreducible and reduced.
\end{proposition}
\begin{proof}
The witnesses give the nonemptiness required by the incidence lemma.
More precisely, the six first-jet rows at their respective tangent
pairs have ranks $6,6,6$, so all three witnesses lie in the main
rank-six pair incidence. This is checked directly in the retained
program \path{agents/monodromy/code/verify_ss_collision_main_incidence.py},
whose output includes the rows and their reduced echelon forms;
it passed in $0.0005378700006986037$ seconds. The reduction-descent
auditor also computed these three ranks independently.
They also ensure that the indicated source curves are smooth at
points of each collision divisor. The source singular-curve
discriminant is irreducible: the singular-point incidence is a
projective bundle over $X$. Its generic member has a unique ordinary
node, by the first/second-jet bounds (the incidence of two singular
points has dimension at most six). Its degree is
$c_2(J^1L)=3L^2=54$: on a generic coefficient pencil its singular
members are the zeros of the rank-three first-jet section on
$X\times\mathbf P^1$, and its top Chern number is $54$.
The ordinary-node Hessian is invertible, so the incidence has generic
separable degree one. In particular this discriminant differs from
both cubic source factors.

Outside the source discriminant and the collision divisors, the
norm curve is nodal with normalization $C_s$. The normalization
sequence has a Frobenius-unit graph term, so the relative determinant
vanishes exactly when $D^2Q$ does. At the generic point of the source
discriminant, no collision divisor occurs, by the smooth-source
witnesses. The source node maps isomorphically to a target node;
all additional gluing terms are again Frobenius-unit. Thus this
discriminant is not an additional factor of $\nu^*h_3$.
There are no other codimension-one exceptions: noninvariance and
smoothness exclude all failures of the finite normalization map,
and nonordinary identifications require either tangency or a triple.
Consequently every irreducible factor of $\nu^*h_3$ is one of
$D,Q,R_g,T_{g,h}$.

We also need the converse inclusion for the collision divisors.
For a reduced local curve singularity with $r$ branches, write
$A\subset\widetilde A$ for its completed local ring and normalization.
The branch constants in $\widetilde A/A$ form $k^r/k$, on which
Frobenius is bijective. Modulo these constants every element has
positive order on each branch; a sufficiently high Frobenius power
lies in the conductor and is therefore zero in the quotient.
Thus the Frobenius-stable rank of $\widetilde A/A$ is exactly $r-1$.
The global normalization sequence, and exactness of the stable-part
functor over the perfect field, now give for a connected reduced
curve with normalization components $C_i$
\[
 g(Z)-f(Z)=\sum_i\bigl(g(C_i)-f(C_i)\bigr)
              +\sum_x(\delta_x-r_x+1).
\]
Here $g=h^1(\mathcal O)$ and $f$ is its Frobenius-stable rank.
At a generic point of any $R_g$ or $T_{g,h}$ the source is smooth,
by the checked witnesses and irreducibility, and it lies outside
$D\cup Q$ by the preceding lemma. Hence its defect is exactly two.
A tangency of distinct smooth branches has
$\delta-r+1\ge1$, and so does an identification of at least three
smooth branches in the smooth target surface. The target curve
therefore has defect at least three. Its relative coherent dimension
subtracts the two-dimensional ambient cohomology, whose stable rank
is zero. Thus its relative Frobenius is not bijective. This proves
that every generic collision image lies in the target Hasse divisor;
no unproved ordinarity assertion about a triple image is used.

Finally consider the local family at the generic point of $D$.
The source curve is smooth and the norm has only transverse double
identifications. After a finite etale base change ordering the nodes
and their branches, the normalization quotient is a direct sum of
copies of the base ring, with Frobenius the ordinary square map.
The resulting graph term in $R^1\mathcal O$ is a finite free
Frobenius-unit subbundle. The ambient $H^1(\mathcal O)$ subbundles
on source and target identify under the odd isogeny. Passing to
relative coherent cohomology therefore gives an exact sequence of
vector bundles whose graph term has invertible Frobenius determinant
and whose quotient is the source relative bundle. Determinants in
this exact sequence multiply. The equality is valid over the local
ring, not just on its generic fibre, so the vanishing order along
$D$ in $\nu^*h_3$ is exactly the order two in $D^2Q$.

The first witness and Lemma~\ref{lem:ss-mixed-self-tacnode} identify
the global target Hasse component containing the image of $Q$ with
the one containing the image of $R_G$. The target divisor is smooth
at that witness, so both assignments are unique. Every source
normalizer automorphism descends under $\phi$ and preserves its
norm line bundle. It preserves $Q$, the unique multiplicity-one
component of the source divisor $D^2Q$, and hence fixes the unique
target component just identified. Transitivity on kernel differences
therefore assigns every $R_g$ to that same target component.
The other two witnesses and the tangent-triple bridge identify the
two triple orbits with it as well. Smoothness at each witness makes
these assignments valid even if a pullback has multiplicity.

Any other target Hasse component would therefore pull back to a
nonzero divisor supported only on $D$, with multiplicity at most two.
Its pullback degree would be at most six. But every nonzero target
hypersurface has pullback degree $9\deg(H)\ge9$, since
$\nu^*\mathcal O(1)=\mathcal O(9)$. This contradiction proves
irreducibility. Smoothness at the first witness proves that the
multiplicity of the single component is one.
\end{proof}

\begin{corollary}[The degree-three Hasse divisor]
\label{cor:ss-degree3-hasse-irreducible}
The relative Hasse divisor on $|L^3|$ is geometrically irreducible
and reduced.
\end{corollary}
\begin{proof}
Lemma~\ref{lem:ss-global-contact03} supplies the first witness, and
Lemma~\ref{lem:ss-global-tangent-triples} supplies the second and
third. All three sources, every global collision, the local
Cartier hypotheses and the nonzero transverse derivatives have
been independently checked. Apply the preceding proposition.
\end{proof}

\begin{remark}[Retained failures and scope]
The first local contact candidate had coefficient vector
$(8,3,45,21,31,9,33,17,4)$, Cartier-value difference $38$ and
contact coefficient $29$. Its later global check found an unwanted
triple. A second contact candidate also failed a global triple
check. Their discovery and verification records were retained;
neither failed candidate is used in the proof. The third contact
candidate and the two tangent-triple candidates are the successful
witnesses cited above.

This corollary determines the geometry of the degree-three Hasse
divisor. The full degree-three group $\mathrm{GL}_{80}(\mathbf F_2)$
was already proved by a separate certificate. The subsequent uniform collision witnesses extend this argument to
every odd degree in Theorems~\ref{thm:ss-all-odd-Hasse} and
\ref{thm:ss-all-odd-full-groups}. The remaining even-degree images are determined in
Theorem~\ref{thm:ss-product-all-full-groups}.
\end{remark}
